\documentclass[a4paper]{scrartcl}

\usepackage{scrhack}
\usepackage{graphicx}
\usepackage[utf8]{inputenc}

\addtokomafont{titlehead}{\flushright}
\addtokomafont{subject}{\vspace{3cm}\flushleft}
\addtokomafont{title}{\flushleft}
\addtokomafont{subtitle}{\flushleft}
\addtokomafont{author}{\flushleft\setlength{\tabcolsep}{0pt}}
\addtokomafont{date}{\flushleft}
\addtokomafont{publishers}{\flushleft}

\titlehead{\includegraphics[scale=2]{../../templates/logo_en}}
\subject{Software Engineering and Design}
\title{Interview Questions}
\subtitle{Mental Health Care Patient Management System (MHC-PMS)}
\author{
\begin{tabular}{l}
\normalfont\bfseries{Team White:}\\
Dellsperger Jan\\
Ellenberger Roger\\
Sheppard David\\
Sidler Matthias\\
Spring Mathias\\
Thöni Stefan
\end{tabular}
}
\date{\today}
\publishers{Version 1.2}

\begin{document}

\begin{titlepage}
	\maketitle
\end{titlepage}

\section{Aufbau und Fragenkatalog}

\subsection{Termine}
Wir haben zwei Interview-Termine. Am 23. März 2016 von 7:00 bis 8:00 und am 24. März 2016 von 9:00 bis 10:00.

\subsection{Ablauf}
Grobe Gliederung des Interviews:
\begin{itemize}
\item Gegenseitige Vorstellung
\item Erlaubnis für Tonaufnahme einholen
\item Erklären des Ablaufs
\item Interview
\end{itemize}


\subsection{Fragen}
\begin{enumerate}
\item Was sind Ihre Tätigkeiten / Verantwortlichkeiten?
\item Was für Software würde Ihnen (als Manager) im Alltag helfen?
\item Sind momentan solche oder ähnliche Systeme im Einsatz? Falls ja, was sind gute und weniger gute Aspekte dieser Systeme? Falls nein, was müsste ein solches System bieten damit es zum Einsatz käme? (Mit welchen Applikationen verwalten Sie Patienten mit psychischen Krankheiten?)
\item Welche Daten sind für Sie von Bedeutung?
\item Welche rechtlichen Verpflichtungen gibt es im Zusammenhang mit Patientendaten? (Was bedeutet das für Datenanalyse / Reporting?) Welche Daten müssen anonym behandelt werden?
\item Gibt es gegenüber Krankenkassen, Bundesämtern, Partnerorganisationen etc. Verpflichtungen Daten abzuliefern? Wenn ja welche und wie werden diese ausgetauscht?
\item Welche Reports / Statistiken verwenden Sie regelmässig?
\item Sehen Sie mögliche Anwendungen für eine mobile Applikation?
\item Sind Sie täglich mit repetitiven Aufgaben konfrontiert?

\item Gibt es Bedarf für Datenaustausch zwischen verschiedenen Benutzern? Falls ja, wie sieht dieser aus?
\item Wer würde eine solche Software benutzen? Welche Positionen halten diese Personen inne?
\item Was sind momentane IT-Trends in der Gesundheitsbranche?

\end{enumerate}

Optionale Frage:
\begin{itemize}
\item Haben Sie spezifische Innovationsprozesse und falls ja, wie sehen diese aus? POCs etc.a
\item Wie sieht Ihr Lifecyclemanagement aus?
\end{itemize}


\section{Transkript}

\subsection{Einwilligung für Tonaufnahme}
Die Interviews wurden aufgezeichnet. Die nötige Erlaubnis der Teilnehmenden wurde vorher mündlich eingeholt. Den Befragten ist bekannt, dass ihr Name und ihre Position verschleiert wird.


\subsection{Anonymisierung}
Die befragten Personen werden nicht mit dem richtigen Namen erwähnt und die Funktion bzw. Erfahrung lediglich grob umschrieben. Der Arbeitsort wird nicht genannt.


\section{Interview 23. März 2016}
Teilnehmer:
\begin{itemize}
\item Konrad Meier [KM]: Mehrere Jahre IT-Erfahrung in diversen Branchen, aktuell in der IT-Leitung eines Spitals
\item Alain Berger [AB]: 15 Jahre tätig in der Psychiatrie, studierter Medizin-Informatiker
\item David Sheppard [DS]: Interviewer
\item Roger Ellenberger [RE]: Interviewer
\end{itemize}

Interview:
\begin{description}
\item[RE 00:00] [...]
\item[RE 00:22] Wir haben den Auftrag der Berner Fachhochschule - von unseren Dozenten im Modul Software-Engineering - ein Softwareprojekt durchzuführen für ein Patientenmanagementsystem für Patienten mit psychischen Problemen. Das Projekt ist unterteilt in zwei Hauptteile. Zum einen einen ist es ein Patientenmanagementsystem mit dem Ärzte, Krankenschwestern und auch Patienten selbst damit arbeiten können. Und auf der anderen Seite gibt es einen Management-Teil. Es gibt mehr Aufgaben im Bereich Patientenverwaltung, aber wir haben den Auftrag rein den Management-Teil auszuarbeiten. Anwendungsfälle der Software auf Seiten Patient ist Terminverwaltung, Medikament- / Rezeptausstellung und Unterstützung des Patienten damit dieser weiss welche Medikamente er einnehmen  muss. Zudem sollen die Mitarbeiter ein Tool erhalten mit dem sie möglichst einfach Termine verwalten können und auch die Patientengeschichte nachvollziehen können.
\item[AB 01:48] Geht es um stationäre oder um ambulante Behandlung?
\item[RE 01:52] Es ist ambulant, nicht alle Patienten sollen stationär behandelt werden.
\item[DS 01:58] Es geht grundsätzlich um beide Behandlungsarten mit Fokus auf ambulante Behandlung.
\item[AB 02:06] Weil üblicherweise haben sie im stationären Bereich bereits elektronische	Patientenakten oder auch papierbasierte Patientenakten, die als Managementsystem gebraucht wird. Dabei wird zusammen mit dem Patienten die Pflegeplanung erstellt. Deshalb habe ich gefragt und nehme an, dass dies nicht betroffen ist. [...]




\end{description}


\end{document}